\section{Background}
\label{sec:background}

\subsection{Socket: The Performance Bottleneck}
\label{subsec:motivation}

%\begin{figure}[t]
%	\centering
%	\includegraphics[width=0.3\textwidth]{images/fixme}
%	\caption{Fraction of CPU time in the kernel (socket connection setup, socket data transmission and non-socket system calls) and user-mode applications.}
%	\label{fig:socket-kernel-time}
%\end{figure}

%\begin{figure}[t]
%	\centering
%	\includegraphics[width=0.3\textwidth]{images/fixme}
%	\caption{Performance of back-end systems using inter-server socket, intra-server socket, inter-server RDMA and intra-server shared memory.}
%	\label{fig:backend-performance}
%\end{figure}


%\begin{figure}[t]
%	\centering
%	\includegraphics[width=0.3\textwidth]{images/fixme}
%	\caption{Latency comparison of mutual exclusion mechanisms (CAS, mutex, futex) and cache migration.}
%	\label{fig:mutual-exclusion}
%\end{figure}

%\begin{figure}[t]
%	\centering
%	\includegraphics[width=0.3\textwidth]{images/fixme}
%	\caption{Throughput (bar) and latency (line) for 16B messages using inter-server TCP socket, inter-server RDMA, intra-server TCP socket, UNIX socket, pipe and shared memory queue.}
%	\label{fig:socket-comparison}
%\end{figure}

\begin{table}[t]
	\centering
	\scalebox{0.9}{
		\begin{tabular}{l|l|l|}
			\hline
			Operation	& Latency  & Throughput  \\
						& ($\mu$s) & (M~op/s) \\
			\hline
			\hline
			Inter-core cache migration	& 0.03 & 50 \\
			\hline
			System call (before KPTI) & 0.05 & 21 \\
			\hline
			CPU L3 cache miss & 0.07 & 14 \\
			\hline
			Atomic operation & 0.10 & 5$\sim$10 \\
			\hline
			Shared memory queue & 0.25 & 27 \\
			\hline
			\textbf{Intra-server \sys} & 0.30 & 22 \\
			\hline
			System call (after KPTI) & 0.20 & 5 \\
			\hline
			Copy one page (4~KiB) & 0.40 & 5 \\
			\hline
			NIC cache miss & 0.45 & 2.2 \\
			\hline
			Cooperative context switch & 0.52 & 2.0 \\
			\hline
			Map 1 page (4~KiB) & 0.78 & 1.3 \\
			\hline
			Map 32 pages (128~KiB) & 1.2 & 0.8 \\
			\hline
			Two-sided RDMA & 1.6 & 8 \\
			\hline
			One-sided RDMA & 1.6 & 13 \\
			\hline
			\textbf{\sys via RDMA} & 1.6 & 8 \\
			\hline
			Semaphore, mutex, futex & 2.8$\sim$5.5 & 0.2$\sim$0.4 \\
			\hline
			Intra-server TCP & 11 & 0.9 \\
			\hline
			Copy 32 pages (128~KiB) & 13 & 0.08 \\
			\hline
			Create TCP socket & 14 & 0.07 \\
			\hline
			Inter-server TCP & 30 & 0.3 \\
			\hline
		\end{tabular}
	}
	\caption{Per-core throughput and round-trip latency of operations on our testbed (settings in Sec.\ref{sec:evaluation}). Message size is 8 bytes if not specified.}
	\vspace{-15pt}
	\label{tab:operation-performance}
\end{table}


\begin{table*}[t!]
	\centering
	\begin{tabular}{ll|ll|ll|ll}
		\hline
		\multicolumn{2}{c|}{Initialization} &
		\multicolumn{2}{c|}{Connection Mgmt} &
		\multicolumn{2}{c|}{Data Transmission} &
		\multicolumn{2}{c}{Process Mgmt} \\
		\hline
		API & Cat. &
		API & Cat. &
		API & Cat. &
		API & Cat. \\
		\hline
		\hline
		\textbf{socket} & Local &
		\textbf{connect} & NoPart &
		\textbf{send(to,(m)msg)} & P2P &
		\textit{(v)fork} & NoPart \\
		\hline
		bind & NoPart &
		\textbf{accept(4)} & P2P &
		\textbf{recv(from,(m)msg)} & P2P &
		\textit{pthread\_create} & NoPart \\
		\hline
		listen & NoPart &
		\textbf{\textit{fcntl, ioctl}} & Local &
		\textbf{\textit{write(v)}} & P2P &
		\textit{clone} & NoPart \\
		\hline
		socketpair & Local &
		\textbf{(get,set)sockopt} & Local &
		\textbf{\textit{read(v)}} & P2P &
		\textit{execve} & NoPart \\
		\hline
		getsockname  & Local &
		\textbf{\textit{close}, shutdown} & P2P &
		\textbf{\textit{memcpy}} & Local &
		\textit{exit} & P2P \\
		\hline
		\textbf{\textit{malloc}} & Local &
		getpeername & Local &
		\textit{(p)select} & P2P &
		\textit{sleep} & P2P \\
		\hline
		\textbf{\textit{realloc}} & Local &
		\textit{dup(2)} & P2P &
		\textit{(p)poll} & P2P &
		\textit{daemon} & P2P \\
		\hline
		\textit{epoll\_create} & Local &
		\textbf{\textit{epoll\_ctl}} & Local &
		\textbf{\textit{epoll\_(p)wait}} & P2P &
		\textit{sigaction} & Local \\
		\hline
	\end{tabular}
	\caption{Linux APIs that are related to socket and intercepted by \libipc{}. Categories include local, peer-to-peer (P2P) and non-partitionable (NoPart). APIs in \textit{italic} indicate usages besides socket. APIs in \textbf{bold} are called relatively more frequently.}
	\vspace{-10pt}
	\label{tab:socket-api}
\end{table*}



Linux kernel socket is a well-known bottleneck for communication intensive applications. We stress test Nginx load balancer, memcached and Redis with one request per TCP connection or a stream of requests on pre-established TCP connections, and find that 50\%$\sim$90\% CPU time is consumed in the kernel, mostly dealing with socket system calls. The one request per connection scenario shows much lower throughput, because of overhead in connection creation.

\textbf{Container networking.}

\textbf{Event-driven communication pattern (code snippet).}


%If we were able to mitigate the overhead associated with sockets, application performance would be 2x$\sim$10x.
%As another example, we replaced socket with RDMA and enabled zero copy gRPC in distributed Tensorflow~\cite{abadi2016tensorflow}, and training a VGG network shows 5x speedup.

%Socket is a well-known bottleneck in fast data center networks \RED{(cite)}. In a data center, we categorize systems that process queries from Internet users to be \textit{front-end} (\textit{e.g.} DNS server, load balancer and Web service), and the other systems (\textit{e.g.} database, stream processing and machine learning) to be \textit{back-end}.

%For front-end systems, as shown in Figure~\ref{fig:socket-kernel-time}, 70\% -- 90\% CPU time is used in socket system calls. They maintain a large number of socket connections to Internet users. Establishing a socket connection takes $\approx$40$\mu$s CPU time, so a CPU core can only serve 25K new connections per second~\cite{lin2016scalable}.

%For back-end systems, the performance using socket is significantly worse than using RDMA or shared memory. This is because back-end systems interact frequently with other nodes and services, consequently the latency is vital for their performance. As shown in Figure~\ref{fig:socket-comparison}, socket latency between different processes in a same server is $\approx$10~$\mu$s, even much higher than inter-server RDMA latency ($\approx$2~$\mu$s).

%\textbf{Event-driven programming.}

%\textbf{Socket process: connection setup and data transmission.}

\subsection{Why Socket is Slow}
\label{subsec:challenges}


\parab{Kernel crossing.}
Socket APIs are implemented in kernel and thus require kernel crossing for each socket operation. Recently, the Kernel Page-Table Isolation patches~\cite{kpti} to protect against Meltdown~\cite{Lipp2018meltdown} attacks make kernel crossings 4x expensive.
% (from 47~ns on Linux 4.11 to 200~ns on Linux 4.16 per system call). 
Together with additional cache invalidations, KPTI patches reduce TCP socket throughput by 40\% for 64-byte messages. These side-channel attacks also make us rethink the Linux security model in which the complicated network stack needs to be trusted and protected. %, while sharing cores with applications.
%To ensure isolation among processes, one category of dataplane operating systems~\cite{belay2017ix,tsai2017lite} use kernel-based network stacks. However, non-batched system calls have context switch overhead, while batched system calls increase latency.
Several works~\cite{martins2014clickos,roghanchi2017ffwd,huang2017high} delegate operations to a virtual switch running on dedicated cores, making the virtual switch a bottleneck.

We aim to bypass kernel with process isolation, while not delegating all operations to a centralized coordinator.


\parab{Context switch and event notification.}
%As shown in Table~\ref{tab:operation-performance}, switching from one process to another cooperatively is 2.5x as expensive as a system call. 
As Table~\ref{tab:operation-performance} shows, inter-process context switch based on semaphore, mutex or \texttt{futex} is 5x$\sim$10x slower than cooperative multitasking. This extra latency is due to event notification mechanism and process scheduler in kernel. %For two processes that run on different cores, TCP socket latency is only a bit higher than mutex. 

We aim to use cooperative multitasking instead of kernel event notification mechanisms.

\parab{Synchronization.}
Linux kernel acquires global locks during connection creation and a per-socket lock for each socket operation~\cite{boyd2010analysis,han2012megapipe,lin2016scalable}. Previous works~\cite{boyd2010analysis,clements2015scalable} suggest that many socket operations are not commutable and thus synchronization cannot always be avoided. For example, synchronization is required when a socket is shared after fork.

We minimize synchronization overhead in two ways. First, we optimize for the common case and remove synchronization from frequent socket operations. Second, we leverage the fact that shared memory message passing is much cheaper than locking, and use message passing as the exclusive synchronization mechanism.

%Spinlock in kernel is implemented with Compare-And-Swap (CAS) instruction, which costs $100\sim200$~ns per acquire and release. In comparison, shared memory message passing only needs one cache migration~\cite{roghanchi2017ffwd} (30~ns).
%In comparison, shared memory message passing
%Scalability analysis of Linux system calls~\cite{boyd2010analysis,clements2015scalable} suggests that many socket operations are not commutable and thus impossible to scale for all cases. 
%For example, socket provides \texttt{SO\_REUSEPORT} option. When this option is enabled, multiple processes on one host can \texttt{bind} to the same port and incoming connections are dispatched to the listeners. During connection setup, coordination is required to allocate file descriptors and port numbers, load balance connections and allocate buffers. In addition, a socket may be shared by multiple processes and threads in a same process share socket connections. When a process forks, the parent and child processes also share the sockets created before fork.


%For inter-server data transmission, Linux kernel copies data four times: on the send side, from application to kernel socket buffer, then to NIC ring buffer; the receive side is similar. For intra-server, data is copied three times: from application to kernel send buffer, then to kernel receive buffer via loopback interface, and finally to another application. Each CPU core can copy $\approx$5~GiB data per second~\cite{panda2016netbricks}, that's why a kernel TCP connection is hard to saturate an 100~Gbps NIC.
%Historically, \texttt{send} was designed as a blocking operation, and the send buffer may be overwritten by the application after \texttt{send} function call.
%In a non-blocking \texttt{send}, the process could overwrite the send buffer after \texttt{send} function call returns. To avoid race condition on the data buffer, network stack needs to copy data to an internal buffer. %before \texttt{send} returns, and then send the copied data in background.
%For the \texttt{recv} operation, application provides a buffer and read the data after \texttt{recv}.
%If we implement \texttt{send} as a lazy operation, \textit{i.e.}, the data is buffered on the sender, the receiver needs to wait for a round-trip time on \texttt{recv}.
%For this reason, most socket implementations send data to the receiver eagerly.
%Because the receiver's network stack cannot predict the user-provided \texttt{recv} address, it needs to buffer received data internally, then copy to the desired destination when \texttt{recv} is called.
%Our goal in \sys is to achieve zero copy for large data transfers without making changes to applications. We accomplish this by taking advantage of the page mapping mechanism in virtual memory. 





\parab{Cache miss.}
The vast number of connections from Internet users require cloud systems to serve millions of connections efficiently~\cite{nishtala2013scaling,lin2016scalable,belay2017ix}. %, often called the C10M problem~\cite{graham2013c10m}. 
In Linux, each socket connection has dedicated send and receive buffers, each at least one page (4~KiB) in size. With millions of concurrent connections, the socket buffers can consume gigabytes of memory, most of which is empty. Accessing many buffers randomly leads to non-cached memory accesses. This problem also exists in RDMA NICs~\cite{mprdma,kaminsky2016design} which have limited on-board memory.
%, which is slower than sequential access to a single buffer~\cite{li2017kv}. 
%This problem is exaggerated in RDMA. RDMA NIC caches connection states in limited on-board memory. Just over a few hundred active RDMA connections could cause cache misses and degrade performance~\cite{mprdma,kaminsky2016design}. 
%Moreover, event-driven applications use \texttt{epoll} to detect which connections are ready for send or receive. If the event queue is separated from data queues, each \texttt{recv} operation involves two cache migrations for event and data~\cite{yasukata2016stackmap}. 

We aim to multiplex socket connections and minimize memory accesses per data transmission.

%\textbf{Virtual file system.}
%Each socket connection is a \textit{file descriptor} (FD) in Linux. First, each new connection involves expensive \texttt{inode} and \texttt{dentry} creation in the \texttt{proc} file system, which is not scalable and rarely used except for diagnostic and monitoring utilities. Second, socket operations need to pass through the virtual file system (VFS), which introduces CPU overhead and additional latency.

\parab{Memory copy.}
Socket \texttt{send} and \texttt{recv} APIs cause memory copies between application and network stack. For \texttt{send}, the application may overwrite the memory right after \texttt{send} returns. Simply removing these copies may violate application correctness. For \texttt{recv}, application provides a buffer and read the data in buffer after \texttt{recv} returns. If we buffer data on the sender and only transmit it on \texttt{recv}, the receiver needs to wait for a round-trip time. For this reason, most socket implementations buffer data on the receiver and copy it to application on \texttt{recv}.


We aim to achieve zero copy for large transfers without changing applications. Rather than copying, we map the physical pages to new virtual addresses. However, as Table~\ref{tab:operation-performance} shows, remapping a single page is more expensive than copying it, because of kernel crossing and TLB flush. Hence, we batch page mapping operations to amortize the overhead.



\parab{Transport.}
Linux TCP has evolved to become quite complicated over the years to address performance and security concerns~\cite{yasukata2016stackmap}. However, intra-server communication may not need many TCP features, \textit{e.g.}, congestion control and loss recovery. Hence, we leverage shared memory queue instead of TCP as intra-server transport.%the reordering, loss recovery and congestion control offered by TCP. %QoS and performance isolation are also not required because synchronization-free inter-core communication in a NUMA node can scale with number of cores~\cite{intel-manual}. 


%For inter-server communication inside , RDMA offers reliable transport.
For inter-server communication inside data centers, RDMA is a good option as it uses hardware offloading to provide lower latency and CPU utilization than software-based transport. RDMA has been deployed in many data centers~\cite{guo2016rdma}. The main challenge of using RDMA is to bridge the semantics of socket and RDMA~\cite{dragojevic2014farm}. For example, RDMA preserves messages boundaries while TCP does not. Further, RDMA verbs have different efficiency and overheads~\cite{kalia2014using,kaminsky2016design}. We aim to use RDMA efficiently for existing socket applications. 


\subsection{Related Work}
\label{sec:related}



\begin{table*}[t]
	\centering
	\scalebox{0.88}{
		\begin{tabular}{l|c|ccc|cc|ccc|}
			\hline
			Category	& \multicolumn{1}{c|}{Linux Opt.} & \multicolumn{3}{c|}{New Kernel Stack} & \multicolumn{2}{c|}{User-space Packet} & \multicolumn{3}{c|}{User-space Socket} \\
			\hline
			System	& FastSocket & IX & MegaPipe & StackMap & Arrakis & FreeFlow & mTCP & libvma & SocksDirect \\
			\hline
			\hline
			Socket-like API & \yes & & \yes & \yes & \yes & & \yes & \yes & \yes \\
			\hline
			Linux Compatible & \yes & & & & & & & \yes & \yes \\
			\hline
			Process Isolation & \yes & \yes & \yes & \yes & \yes & \yes & \yes & \yes & \yes \\
			\hline
			\hline
			Kernel Bypass & & & & & \yes & \yes & \yes & \yes & \yes \\
			\hline
			Cooperative Multitasking & & & & & & & & & \yes \\
			\hline
			Zero Copy & & \yes & & \yes & \yes & \yes & & & \yes \\
			\hline
			\hline
			NIC-bypass IPC & \yes & & \yes & \yes & & \yes & \yes & & \yes \\
			\hline
			RDMA as Transport & & & & & & \yes & & & \yes \\
			\hline
			\hline
			Scalable Socket Creation & \yes & \yes & \yes & \yes & \yes & \yes & \yes & & \yes \\
			\hline
			Lock-free Multi-thread & & & & & & & & & \yes \\
			\hline
			Scale to Many Sockets & \yes & \yes & \yes & \yes & & & \yes & & \yes \\
			\hline
		\end{tabular}
	}
	\caption{Comparison of high performance socket systems.}
	\label{tab:related-work}
	\vspace{-10pt}
\end{table*}


There has been extensive work aiming to make the best use of multi-core CPU and data center network. Table~\ref{tab:related-work} categorizes several representative approaches.

\parab{Linux kernel optimization.}
One line of research optimizes the kernel stack for higher socket performance. FastSocket~\cite{lin2016scalable} and Affinity-Accept~\cite{pesterev2012improving} scale connection creation to multiple cores, but synchronization is still needed when multiple threads share a socket.
FlexSC~\cite{soares2010flexsc} proposes exception-less system calls to reduce kernel crossing overhead.
Zero-copy socket~\cite{thadani1995efficient,chu1996zero} still needs copy-on-write on senders.
In addition, they fail to remove cache miss and transport overheads.


\parab{New OS stacks.}
Another line of research proposes new OS stacks with modified socket interface, mostly aiming at zero copy and fast event notification. Existing socket applications need modifications to use the new interfaces.
For intra-server connections, Arrakis~\cite{peter2016arrakis} and IX~\cite{belay2017ix} use the NIC to forward packets from one process to another. The hairpin latency from CPU to NIC is at least two PCIe delays, which is one order of magnitude higher than inter-core cache migration delay. In addition, the data plane switching throughput of a NIC is constrained by PCIe bandwidth (Figure~\ref{fig:eval-corenum-tput}).

For inter-server connections, most OS stacks implement transport in software. IX~\cite{belay2017ix} and Stackmap~\cite{yasukata2016stackmap} run in the kernel to enforce QoS policy or preserve protocol compatibility with Linux, while Arrakis~\cite{peter2016arrakis} and SandStorm~\cite{marinos2014network} run in user mode to maximize performance.
RDMA and transport virtualization~\cite{tsai2017lite,niu2017network} also enforce QoS in the hypervisor.
Due to the additional level of indirection, kernel stacks cannot remove kernel crossing, while batched syscalls add latency.
Further, large-scale deployment of kernel-based stacks is more complicated than user-space libraries~\cite{andromeda}.
\sys offloads transport and QoS to NIC hardware.
RDMA transport has been deployed in many data centers~\cite{guo2016rdma}, and an emerging line of work~\cite{zhu2015congestion,lu2017memory,mprdma} improves congestion control and QoS in large-scale RDMA deployments.
For flexibility, programmable NICs are being adopted in data centers~\cite{smartnic,cavium}, as they are more efficient than general-purpose CPUs for network processing~\cite{kaufmann2015flexnic,li2016clicknp}.



\parab{User-space socket.}
A third line of research runs socket in user space.
mTCP~\cite{jeong2014mtcp}, Seastar~\cite{seastar}, 
F-stack~\cite{fstack} and LOS~\cite{huang2017high} use a high performance packet I/O framework (\textit{e.g.} netmap~\cite{rizzo2012netmap}, DPDK~\cite{dpdk} and PF\_RING~\cite{pf-ring}) and achieves compatibility with most Linux socket functions and scalability with number of cores and sockets.
LibVMA~\cite{libvma}, OpenOnload~\cite{openonload} and DBL~\cite{dbl} are fully compatible with existing applications. However, they use vendor-specific NIC features and do not scale to multiple threads or connections.
In addition, user-space sockets do not support zero copy or efficient multitasking.

Most user-space sockets focus on inter-server and do not optimize for intra-server connections.
FreeFlow~\cite{freeflow} uses shared memory for intra-server communication and RDMA for inter-server, but it provides an RDMA interface.
Existing socket to RDMA translation approaches, \textit{e.g.} SDP~\cite{socketsdirect} and rsockets~\cite{rsockets} are not fully compatible with Linux and do not address scalability challenges.

